\documentclass[11pt,letterpaper,openany]{book}

\usepackage[margin=1.2in]{geometry}
\usepackage{amsmath, amssymb, amsthm}
\usepackage{mathtools}
\usepackage{bm}
\usepackage{tikz}
\usepackage{pgfplots}
\pgfplotsset{compat=1.18}
\usetikzlibrary{arrows.meta, decorations.markings, calc, positioning}
\usepackage{tcolorbox}
\tcbuselibrary{theorems, breakable, skins}
\usepackage{hyperref}
\usepackage{imakeidx}
\makeindex[intoc, title=Index]
\usepackage{enumitem}
\usepackage{fancyhdr}
\usepackage{titlesec}
\usepackage{epigraph}
\usepackage{xcolor}
\usepackage{booktabs}
\usepackage{listings}

% Code formatting
\definecolor{codegreen}{rgb}{0,0.6,0}
\definecolor{codegray}{rgb}{0.5,0.5,0.5}
\definecolor{codepurple}{rgb}{0.58,0,0.82}
\definecolor{backcolour}{rgb}{0.95,0.95,0.92}
\lstdefinestyle{mystyle}{
    backgroundcolor=\color{backcolour},
    commentstyle=\color{codegreen},
    keywordstyle=\color{magenta},
    numberstyle=\tiny\color{codegray},
    stringstyle=\color{codepurple},
    basicstyle=\ttfamily\footnotesize,
    breakatwhitespace=false, breaklines=true,
    captionpos=b, keepspaces=true, numbers=left,
    numbersep=5pt, showspaces=false, showstringspaces=false,
    showtabs=false, tabsize=2
}
\lstset{style=mystyle}

% Colors
\definecolor{chapblue}{RGB}{0, 70, 130}
\definecolor{accentred}{RGB}{180, 40, 40}
\definecolor{softgray}{RGB}{245, 245, 245}
\definecolor{trajgreen}{RGB}{30, 130, 60}
\definecolor{toolorange}{RGB}{200, 120, 20}

% Theorem environments
\newtcbtheorem[number within=chapter]{handson}{Hands-On}{
  colback=toolorange!5, colframe=toolorange!80,
  fonttitle=\bfseries, separator sign={.~}, breakable
}{ho}
\newtcbtheorem[number within=chapter]{example}{Example}{
  colback=softgray, colframe=black!40,
  fonttitle=\bfseries, separator sign={.~}, breakable
}{ex}
\newtcbtheorem[number within=chapter]{definition}{Definition}{
  colback=chapblue!5, colframe=chapblue!75,
  fonttitle=\bfseries, separator sign={.~}, breakable
}{def}

\newtcbtheorem[number within=chapter]{laymansbox}{In Plain Language}{
  colback=green!5!white, colframe=green!50!black,
  fonttitle=\bfseries, separator sign={.~},
  breakable
}{laymans}

\theoremstyle{plain}
\newtheorem{proposition}{Proposition}[chapter]

% Custom commands
\newcommand{\state}{\bm{x}}
\newcommand{\control}{\bm{u}}
\newcommand{\traj}{\bm{\gamma}}
\newcommand{\Reals}{\mathbb{R}}
\newcommand{\dd}{\mathrm{d}}
\newcommand{\dt}{\,\dd t}
\newcommand{\norm}[1]{\left\lVert #1 \right\rVert}

% Chapter formatting
\titleformat{\chapter}[display]
  {\normalfont\huge\bfseries\color{chapblue}}
  {\chaptertitlename\ \thechapter}{20pt}{\Huge}
\titlespacing*{\chapter}{0pt}{-20pt}{30pt}

% Header/Footer
\pagestyle{fancy}
\fancyhf{}
\fancyhead[LE,RO]{\thepage}
\fancyhead[RE]{\textit{Volume V: Practical Application}}
\fancyhead[LO]{\textit{\nouppercase{\leftmark}}}
\renewcommand{\headrulewidth}{0.4pt}
\setlength{\headheight}{14.5pt}

% Title
\title{%
  {\Huge\bfseries\color{chapblue} The Geometry of Motion}\\[0.4cm]
  {\Large Volume V: Practical Application\\
  The UpstreamDrift Educational Platform}\\[0.6cm]
  {\color{chapblue}\rule{\textwidth}{2pt}}
}
\date{}

\hypersetup{colorlinks=true, linkcolor=chapblue, filecolor=magenta, urlcolor=chapblue, citecolor=accentred}

\newtcbtheorem[number within=chapter]{laymansbox}{In Laymans Terms}{colback=blue!5,colframe=blue!75,fonttitle=fseries,breakable}{layman}
\begin{document}

\maketitle
\thispagestyle{empty}
\cleardoublepage

\frontmatter
\chapter*{Preface}
\addcontentsline{toc}{chapter}{Preface}

This is a book about \emph{doing}. Volumes 0--IV developed the theory of motion: from
differential geometry through nonlinear control, biomechanics, and motor control. This
volume puts every tool to work.

The UpstreamDrift platform provides a unified simulation environment for multi-body
dynamics, trajectory optimization, and controller design. It wraps multiple physics
engines (MuJoCo, Drake, Pinocchio, OpenSim) behind a consistent Python API, allowing
the reader to run every example from the preceding four volumes and to build their own
analyses.

Every chapter in this volume produces a concrete, runnable result. There are no
hypothetical exercises---every example uses real code, real models, and real data.
By the end, you will have implemented a complete pipeline from model construction
through trajectory optimization to robustness analysis, applied to the golf swing
as the running case study.

\tableofcontents
\chapter*{Nomenclature}
\addcontentsline{toc}{chapter}{Nomenclature}

\section*{General Conventions}
\begin{itemize}
    \item Scalars are lowercase or Greek letters (e.g., $m, t, \theta$).
    \item Vectors are bold lowercase (e.g., $\bx, \bu, \bq$).
    \item Matrices are bold uppercase (e.g., $\bA, \bB, \bM, \bJ$).
    \item Expected values and sets are denoted by blackboard bold or script (e.g., $\mathbb{E}, \mathcal{S}$).
\end{itemize}

\section*{State and Kinematics}
\begin{description}
    \item[$\bq \in \Reals^n$] Joint configuration vector (generalized coordinates).
    \item[$\dot{\bq} \in \Reals^n$] Joint velocity vector (generalized velocities).
    \item[$\ddot{\bq} \in \Reals^n$] Joint acceleration vector.
    \item[$\bx \in \Reals^{nx}$] System state vector; commonly $\bx = [\bq^T, \dot{\bq}^T]^T$ for mechanical systems.
    \item[$\bu \in \Reals^m$] Control input vector (e.g., joint torques, muscle activations).
    \item[$\btau \in \Reals^n$] Generalized forces/torques acting on the joints.
    \item[$t \in \Reals$] Time.
\end{description}

\section*{Inertia and Dynamics}
\begin{description}
    \item[$\bM(\bq)$] Mass (inertia) matrix, symmetric positive definite.
    \item[$\bC(\bq, \dot{\bq})$] Coriolis and centrifugal force matrix.
    \item[$\bg(\bq)$] Gravity vector.
    \item[$\bJ(\bq)$] Jacobian matrix relating joint velocities to task-space velocities ($\dot{\bp} = \bJ(\bq)\dot{\bq}$).
\end{description}

\section*{Control and Optimization}
\begin{description}
    \item[$\bA_k = \frac{\partial f}{\partial \bx}$] Local Jacobian matrix of the state dynamics.
    \item[$\bB_k = \frac{\partial f}{\partial \bu}$] Local Jacobian matrix of the control input.
    \item[$V(\bx)$ or $\mathcal{V}$] Value function or Lyapunov function.
    \item[$\bQ$] State penalty matrix in LQR/DDP.
    \item[$\bR$] Control penalty matrix in LQR/DDP.
    \item[$\bK$] Feedback gain matrix ($\delta\bu = -\bK\delta\bx$).
    \item[$\traj$] A trajectory or flow, mapping time and initial conditions to states.
\end{description}

\section*{Biomechanics and Motor Control}
\begin{description}
    \item[$a_i \in [0, 1]$] Muscle activation level.
    \item[$\ell_i^M$] Muscle fiber length.
    \item[$v_i^M$] Muscle fiber contraction velocity.
    \item[$\ell_i^{MT}$] Musculotendon unit length.
    \item[$\bm{F}_{\text{muscle}}$] Force generated by muscles.
    \item[$\bW$] Muscle synergy matrix (weights).
    \item[$\bc(t)$] Muscle synergy activation vector over time.
\end{description}

\section*{Geometry and Manifolds}
\begin{description}
    \item[$SE(3)$] Special Euclidean group of rigid body motions.
    \item[$SO(3)$] Special Orthogonal group of 3D rotations.
    \item[$\mathfrak{se}(3)$] Lie algebra of $SE(3)$, space of spatial twists (velocities).
    \item[$\mathcal{M}$] A smooth manifold.
    \item[$T_{\bx}\mathcal{M}$] Tangent space at point $\bx$.
\end{description}

\cleardoublepage

\mainmatter

\chapter{The UpstreamDrift Platform}
\label{ch:platform}

\epigraph{%
  The best theory is one you can run.
}{}

\section{Architecture Overview}

UpstreamDrift is a unified platform for multi-body dynamics simulation, trajectory
optimization, and motor control analysis. It wraps multiple physics engines behind
a consistent Python API.

\subsection{Core Components}

\begin{center}
\renewcommand{\arraystretch}{1.3}
\begin{tabular}{@{}lll@{}}
\toprule
\textbf{Component} & \textbf{Location} & \textbf{Purpose} \\
\midrule
Physics Engines & \texttt{src/engines/physics\_engines/} & Multi-body simulation \\
Shared Infrastructure & \texttt{src/engines/common/} & State, export, diagnostics \\
Pendulum Models & \texttt{src/engines/pendulum\_models/} & Simplified test cases \\
Learning Modules & \texttt{src/learning/} & ML integration \\
Reinforcement Learning & \texttt{src/reinforcement\_learning/} & RL benchmarks \\
Robotics & \texttt{src/robotics/} & Robot control modules \\
API & \texttt{src/api/} & External access layer \\
Tools & \texttt{src/tools/} & Development utilities \\
\bottomrule
\end{tabular}
\end{center}

\subsection{Supported Physics Engines}

\begin{enumerate}
  \item \textbf{MuJoCo} (\texttt{physics\_engines/mujoco/}): Google DeepMind's
        multi-joint dynamics engine. Soft contacts, muscle actuators, fast simulation.
        Primary engine for biomechanical modeling.

  \item \textbf{Drake} (\texttt{physics\_engines/drake/}): MIT's multibody dynamics
        toolbox. Mathematical rigor, analytical gradients, trajectory optimization.
        Primary engine for control design.

  \item \textbf{Pinocchio} (\texttt{physics\_engines/pinocchio/}): LAAS-CNRS C++ library
        with Python bindings. Efficient rigid-body dynamics with analytical derivatives.
        Primary engine for real-time control.

  \item \textbf{OpenSim} (\texttt{physics\_engines/opensim/}): Stanford's musculoskeletal
        modeling software. Hill-type muscle models, inverse dynamics. Primary engine
        for biomechanics.

  \item \textbf{MyoSuite} (\texttt{physics\_engines/myosuite/}): Muscle-driven
        reinforcement learning environments. Primary engine for motor learning research.
\end{enumerate}

\section{Getting Started}

\begin{lstlisting}[language=Python, caption={Basic UpstreamDrift usage pattern.}]
"""UpstreamDrift: Getting started with multi-engine simulation.

This script demonstrates the basic workflow:
1. Load a model
2. Configure simulation parameters
3. Run a simulation
4. Extract and analyze results
"""
import numpy as np
from pathlib import Path

# Configuration
MODEL_PATH = Path("src/engines/pendulum_models")
ENGINE = "mujoco"  # or "drake", "pinocchio"

def create_simple_pendulum(
    length: float = 1.0,
    mass: float = 1.0,
    damping: float = 0.1
) -> dict:
    """Create a simple pendulum model configuration.

    Parameters
    ----------
    length : float
        Pendulum length in meters.
    mass : float
        Bob mass in kilograms.
    damping : float
        Joint damping coefficient.

    Returns
    -------
    dict
        Model configuration dictionary.
    """
    return {
        "name": "simple_pendulum",
        "n_dof": 1,
        "segments": [{
            "name": "bob",
            "mass": mass,
            "length": length,
            "com_position": length / 2,
            "inertia": mass * length**2 / 3,
        }],
        "joints": [{
            "name": "pivot",
            "type": "revolute",
            "axis": [0, 0, 1],
            "damping": damping,
        }],
        "gravity": [0, -9.81, 0],
    }


def simulate_pendulum(
    config: dict,
    q0: float = np.pi / 4,
    t_end: float = 5.0,
    dt: float = 0.001
) -> dict:
    """Simulate a simple pendulum using Euler integration.

    Parameters
    ----------
    config : dict
        Model configuration.
    q0 : float
        Initial angle (radians from vertical).
    t_end : float
        Simulation duration (seconds).
    dt : float
        Time step.

    Returns
    -------
    dict
        Simulation results with keys 't', 'q', 'qd', 'energy'.
    """
    seg = config["segments"][0]
    m = seg["mass"]
    L = seg["length"]
    b = config["joints"][0]["damping"]
    g = abs(config["gravity"][1])
    I = seg["inertia"]

    n_steps = int(t_end / dt)
    t = np.zeros(n_steps)
    q = np.zeros(n_steps)
    qd = np.zeros(n_steps)
    energy = np.zeros(n_steps)

    q[0] = q0
    qd[0] = 0.0

    for k in range(n_steps - 1):
        # Equations of motion: I * qdd + b * qd + m*g*L/2*sin(q) = 0
        qdd = (-b * qd[k] - m * g * L / 2 * np.sin(q[k])) / I

        # Semi-implicit Euler
        qd[k + 1] = qd[k] + qdd * dt
        q[k + 1] = q[k] + qd[k + 1] * dt
        t[k + 1] = t[k] + dt

        # Energy
        KE = 0.5 * I * qd[k + 1]**2
        PE = m * g * L / 2 * (1 - np.cos(q[k + 1]))
        energy[k + 1] = KE + PE

    return {"t": t, "q": q, "qd": qd, "energy": energy}


# Run simulation
config = create_simple_pendulum(length=1.0, mass=1.0, damping=0.05)
results = simulate_pendulum(config, q0=np.pi/3, t_end=10.0)

print(f"Simulation complete: {len(results['t'])} timesteps")
print(f"Initial energy: {results['energy'][0]:.4f} J")
print(f"Final energy:   {results['energy'][-1]:.4f} J")
print(f"Energy loss:    {(1 - results['energy'][-1]/results['energy'][0])*100:.1f}%")
\end{lstlisting}

\section{Engine Selection Guide}

\begin{handson}{Choosing the Right Engine}{engine-choice}
  \begin{center}
  \renewcommand{\arraystretch}{1.3}
  \begin{tabular}{@{}lll@{}}
  \toprule
  \textbf{If you need\ldots} & \textbf{Use} & \textbf{Why} \\
  \midrule
  Fast contact simulation & MuJoCo & Soft contact model, GPU \\
  Trajectory optimization & Drake & Analytical gradients \\
  Real-time forward dynamics & Pinocchio & C++ speed, O(n) \\
  Muscle-driven biomechanics & OpenSim & Hill models built-in \\
  RL with muscles & MyoSuite & Gym-compatible \\
  Quick prototyping & Built-in Python & No dependencies \\
  \bottomrule
  \end{tabular}
  \end{center}
\end{handson}

\section*{Exercises}

\begin{enumerate}
  \item \textbf{First Simulation.} Run the pendulum simulation and plot angle vs.\ time
        and energy vs.\ time. Verify that energy decreases due to damping.

  \item \textbf{Integration Method.} Replace the semi-implicit Euler with RK4. Compare
        energy conservation over 100 seconds.

  \item \textbf{Double Pendulum.} Extend the code to a double pendulum. Plot the
        chaotic trajectory in configuration space.

  \item \textbf{Engine Comparison.} Run the same pendulum simulation in MuJoCo and
        in the built-in Python code. Compare results quantitatively.
\end{enumerate}

\chapter{Physics Engine Comparison}
\label{ch:engine-comparison}

\epigraph{All models are wrong. Some are useful. The useful ones depend on the question.}{}

\section{Benchmark: Simple Pendulum}

We simulate the same pendulum in all engines and compare:

\begin{lstlisting}[language=Python, caption={Cross-engine benchmark framework.}]
import numpy as np
from dataclasses import dataclass
from typing import Protocol

@dataclass
class BenchmarkResult:
    """Results from a physics engine benchmark."""
    engine_name: str
    positions: np.ndarray    # (n_steps, n_dof)
    velocities: np.ndarray   # (n_steps, n_dof)
    times: np.ndarray        # (n_steps,)
    wall_time_seconds: float
    energy_drift: float      # relative energy change

class PhysicsEngine(Protocol):
    def simulate(self, q0: np.ndarray, qd0: np.ndarray,
                 t_end: float, dt: float) -> BenchmarkResult: ...

def compare_engines(
    engines: dict[str, PhysicsEngine],
    q0: np.ndarray,
    t_end: float = 5.0,
    dt: float = 0.001
) -> dict[str, BenchmarkResult]:
    """Run the same simulation across multiple engines.

    Parameters
    ----------
    engines : dict
        Mapping of engine name to engine instance.
    q0 : np.ndarray
        Initial configuration.
    t_end : float
        Simulation duration.
    dt : float
        Time step.

    Returns
    -------
    dict
        Mapping of engine name to BenchmarkResult.
    """
    results = {}
    qd0 = np.zeros_like(q0)

    for name, engine in engines.items():
        result = engine.simulate(q0, qd0, t_end, dt)
        results[name] = result
        print(f"{name:15s}: wall_time={result.wall_time_seconds:.3f}s, "
              f"energy_drift={result.energy_drift:.2e}")

    return results
\end{lstlisting}

\section{Contact Model Comparison}

\begin{center}
\renewcommand{\arraystretch}{1.3}
\begin{tabular}{@{}llll@{}}
\toprule
\textbf{Engine} & \textbf{Contact Model} & \textbf{Strengths} & \textbf{Limitations} \\
\midrule
MuJoCo & Soft contact (spring-damper) & Fast, stable & Non-physical penetration \\
Drake & Rigid + hydroelastic & Physics-based & Slower \\
Pinocchio & No built-in contacts & Pure dynamics & Needs external contact \\
OpenSim & Hunt-Crossley & Biomechanical & Limited geometries \\
\bottomrule
\end{tabular}
\end{center}

\section*{Exercises}
\begin{enumerate}
  \item Run the benchmark for a double pendulum across all available engines.
  \item Compare energy conservation across engines with different time steps.
  \item Implement a bouncing ball benchmark to compare contact models.
\end{enumerate}

\chapter{Building a Multi-Body Model}
\label{ch:building-models}

\epigraph{The model is the hypothesis made physical.}{}

\section{Model Formats}

\subsection{URDF (Universal Robot Description Format)}
XML-based format used by ROS, Drake, and Pinocchio. Defines links, joints, and visual/collision geometry.

\subsection{MJCF (MuJoCo XML)}
MuJoCo's native format. Extends URDF with actuators, tendons, muscles, and contact parameters.

\section{Building a 7-Segment Golf Model}

\begin{lstlisting}[language=Python, caption={Programmatic multi-body model construction.}]
import numpy as np
from dataclasses import dataclass, field

@dataclass
class Link:
    name: str
    mass: float
    length: float
    com_position: float
    inertia: np.ndarray = field(default_factory=lambda: np.eye(3))

@dataclass
class Joint:
    name: str
    parent: str
    child: str
    joint_type: str  # revolute, ball, fixed
    axis: list[float] = field(default_factory=lambda: [0, 0, 1])
    limits: tuple[float, float] = (-np.pi, np.pi)

@dataclass
class MultiBodyModel:
    name: str
    links: list[Link] = field(default_factory=list)
    joints: list[Joint] = field(default_factory=list)

    @property
    def n_dof(self) -> int:
        dof_map = {"revolute": 1, "ball": 3, "fixed": 0}
        return sum(dof_map.get(j.joint_type, 1) for j in self.joints)

    def to_urdf(self) -> str:
        """Export model as URDF XML string."""
        lines = [f'<robot name="{self.name}">']
        for link in self.links:
            lines.append(f'  <link name="{link.name}">')
            lines.append(f'    <inertial>')
            lines.append(f'      <mass value="{link.mass}"/>')
            lines.append(f'    </inertial>')
            lines.append(f'  </link>')
        for joint in self.joints:
            lines.append(f'  <joint name="{joint.name}" '
                         f'type="{joint.joint_type}">')
            lines.append(f'    <parent link="{joint.parent}"/>')
            lines.append(f'    <child link="{joint.child}"/>')
            lines.append(f'  </joint>')
        lines.append('</robot>')
        return '\n'.join(lines)


def build_golf_swing_model() -> MultiBodyModel:
    """Build a 7-segment golf swing model.

    Segments: pelvis, trunk, upper_arm, forearm, hand, club_shaft, club_head
    """
    model = MultiBodyModel(name="golf_swing_7seg")

    # Segment properties (de Leva 1996, 80kg male)
    model.links = [
        Link("pelvis", mass=11.2, length=0.15, com_position=0.075),
        Link("trunk", mass=25.5, length=0.48, com_position=0.22),
        Link("upper_arm", mass=2.16, length=0.28, com_position=0.16),
        Link("forearm", mass=1.30, length=0.27, com_position=0.12),
        Link("hand", mass=0.49, length=0.19, com_position=0.07),
        Link("club_shaft", mass=0.10, length=1.05, com_position=0.50),
        Link("club_head", mass=0.20, length=0.10, com_position=0.05),
    ]

    model.joints = [
        Joint("pelvis_rotation", "world", "pelvis", "revolute", [0, 1, 0]),
        Joint("trunk_flexion", "pelvis", "trunk", "revolute", [0, 0, 1]),
        Joint("shoulder", "trunk", "upper_arm", "ball"),
        Joint("elbow", "upper_arm", "forearm", "revolute", [0, 0, 1]),
        Joint("wrist", "forearm", "hand", "revolute", [0, 0, 1]),
        Joint("grip", "hand", "club_shaft", "fixed"),
    ]

    return model


model = build_golf_swing_model()
print(f"Golf swing model: {len(model.links)} segments, "
      f"{model.n_dof} DOF")
print(f"\nURDF preview:\n{model.to_urdf()[:500]}...")
\end{lstlisting}

\section*{Exercises}
\begin{enumerate}
  \item Build a 3-segment arm model and export as URDF.
  \item Add muscle attachment points to the golf model.
  \item Load the URDF into MuJoCo and verify the model visually.
\end{enumerate}

\chapter{Forward and Inverse Simulation}
\label{ch:simulation}

\epigraph{Forward simulation answers ``what happens if\ldots?''\\
Inverse dynamics answers ``what caused this?''}{}

The forward dynamics problem computes accelerations from forces, while inverse dynamics
(also called the inverse kinematics problem when phrased in terms of finding required
torques) computes forces from desired accelerations. Both are fundamental to control design
and trajectory planning.

\section{Forward Dynamics with the Companion Platform}

Forward dynamics integration typically uses explicit or implicit Runge-Kutta methods,
or specialized algorithms such as the Recursive Newton-Euler Algorithm (RNEA) for
rigid-body systems \cite{Featherstone2008, Featherstone1983}.

\begin{lstlisting}[language=Python, caption={Forward dynamics simulation pipeline.}]
import numpy as np
from scipy.integrate import solve_ivp

def forward_dynamics_rnea(
    q: np.ndarray,
    qd: np.ndarray,
    tau: np.ndarray,
    model_params: dict
) -> np.ndarray:
    """Compute joint accelerations using RNEA-based forward dynamics.

    Parameters
    ----------
    q : np.ndarray, shape (n,)
        Joint positions.
    qd : np.ndarray, shape (n,)
        Joint velocities.
    tau : np.ndarray, shape (n,)
        Applied joint torques.
    model_params : dict
        Model parameters (masses, lengths, inertias).

    Returns
    -------
    np.ndarray, shape (n,)
        Joint accelerations.
    """
    n = len(q)
    masses = model_params['masses']
    lengths = model_params['lengths']
    g = model_params.get('gravity', 9.81)

    # Simple 2-link model for demonstration
    if n == 2:
        m1, m2 = masses[:2]
        l1, l2 = lengths[:2]
        lc1, lc2 = l1/2, l2/2
        I1 = m1 * l1**2 / 12
        I2 = m2 * l2**2 / 12

        # Mass matrix
        M = np.array([
            [I1 + I2 + m1*lc1**2 + m2*(l1**2 + lc2**2 + 2*l1*lc2*np.cos(q[1])),
             I2 + m2*(lc2**2 + l1*lc2*np.cos(q[1]))],
            [I2 + m2*(lc2**2 + l1*lc2*np.cos(q[1])),
             I2 + m2*lc2**2]
        ])

        # Coriolis
        h = m2 * l1 * lc2 * np.sin(q[1])
        C = np.array([[-h*qd[1], -h*(qd[0]+qd[1])],
                       [h*qd[0], 0]])

        # Gravity
        gvec = np.array([
            (m1*lc1 + m2*l1)*g*np.sin(q[0]) + m2*lc2*g*np.sin(q[0]+q[1]),
            m2*lc2*g*np.sin(q[0]+q[1])
        ])

        qdd = np.linalg.solve(M, tau - C @ qd - gvec)
    else:
        # Placeholder for general n-DOF
        qdd = np.zeros(n)

    return qdd

# Example simulation
params = {'masses': [2.0, 1.0], 'lengths': [0.3, 0.25], 'gravity': 9.81}
q0 = np.array([np.pi/4, np.pi/6])
qd0 = np.zeros(2)
qdd = forward_dynamics_rnea(q0, qd0, np.zeros(2), params)
print(f"Gravity-driven accelerations: {qdd}")
\end{lstlisting}

\section{Simulation Loop and Integration Schemes}

Figure~\ref{fig:sim-loop} illustrates the typical forward simulation loop. At each
timestep, we compute accelerations from the equations of motion, integrate to update
velocities and positions, and advance to the next timestep.

\begin{figure}[h]
  \centering
  \begin{tikzpicture}[
    every node/.style={align=center, draw, rounded corners, minimum width=2.2cm, text centered, font=\small},
    ]
    % State
    \node[fill=blue!20] (state) at (0, 0) {Current State\\$\bm{q}, \dot{\bm{q}}$};
    % Compute accelerations
    \node[fill=green!20] (forces) at (3, 0) {Compute Forces\\$\bm{f}(\bm{q}, \dot{\bm{q}}, t)$};
    % Solve dynamics
    \node[fill=yellow!20] (dynamics) at (6, 0) {Forward Dynamics\\$\ddot{\bm{q}} = M^{-1}(\bm{f} - C - G)$};
    % Integrate
    \node[fill=orange!20] (integrate) at (9, 0) {Integrate\\(RK4, Euler, etc)};
    % Next state
    \node[fill=blue!20] (next) at (12, 0) {Next State\\$\bm{q}_{k+1}, \dot{\bm{q}}_{k+1}$};

    % Arrows
    \draw[->] (state) -- (forces);
    \draw[->] (forces) -- (dynamics);
    \draw[->] (dynamics) -- (integrate);
    \draw[->] (integrate) -- (next);
    \draw[->] (12.8, -0.3) -- (12.8, -1.2) -- (0, -1.2) -- (0, -0.4) node[align=center, above, pos=0.5, fill=white] {repeat};
  \end{tikzpicture}
  \caption{Forward simulation loop: from current state to forces (gravity, contacts, actuators),
    compute accelerations via forward dynamics, integrate, and update to next state.}
  \label{fig:sim-loop}
\end{figure}

\section{Inverse Dynamics Pipeline}

The inverse dynamics problem solves for torques given desired motion: given $\bm{q}(t),
\dot{\bm{q}}(t), \ddot{\bm{q}}(t)$, compute $\bm{\tau}(t)$. This is implemented via the
Recursive Newton-Euler Algorithm (RNEA) \cite{Featherstone2008, luh1980}. The RNEA
can be run either forward (for forward dynamics) or backward (for inverse dynamics)
through the kinematic chain.

\section*{Exercises}
\begin{enumerate}
  \item Run forward dynamics for a 2-link pendulum and plot trajectories.
  \item Implement inverse dynamics and verify $\text{ID}(\text{FD}(\tau)) = \tau$.
  \item Compare with Pinocchio's RNEA implementation for timing.
\end{enumerate}

\chapter{Trajectory Optimization in Practice}
\label{ch:traj-opt}

\epigraph{Theory without practice is empty. Practice without theory is blind.}{}

\section{DDP/iLQR Implementation}

Differential Dynamic Programming (DDP) and iterative Linear Quadratic Regulator (iLQR) are
related algorithms for trajectory optimization \cite{Mayne1966, jacobson1970dpp}. They solve
the finite-horizon optimal control problem by iteratively linearizing around a nominal
trajectory and solving quadratic subproblems. Figure~\ref{fig:ilqr-loop} shows the algorithm structure.

iLQR is a numerical method for implementing the feedback synthesis results from Agrachev \&
Sachkov \cite{AgrachevSachkov2004}, Chapter~6: it iteratively computes the optimal feedback law
via local linearization and quadratic approximation of the cost. The backward-pass computation
of gains $K(t)$ corresponds to the solution of the Riccati equation, a central object in their
optimal control theory.

\begin{figure}[h]
  \centering
  \begin{tikzpicture}[
    every node/.style={align=center, draw, rounded corners, minimum width=2cm, text centered, font=\small},
    ]
    % Initial trajectory
    \node[fill=blue!20] (init) at (0, 4) {Initialize\\$\bm{u}_0$};

    % Forward pass
    \node[fill=green!20] (forward) at (0, 2) {Forward Rollout:\\$\bm{x}(t)$};

    % Backward pass
    \node[fill=yellow!20] (backward) at (4, 2) {Backward Pass:\\Compute $k, K$ gains};

    % Line search
    \node[fill=orange!20] (linesearch) at (2, 0.3) {Line Search:\\Update $\bm{u}$};

    % Convergence check
    \node[fill=pink!20] (converge) at (2, -1.2) {Converged?};

    % Exit
    \node[fill=green!20] (exit) at (4, -2.5) {Optimal\\$\bm{x}^*, \bm{u}^*$};

    % Arrows
    \draw[->] (init) -- (forward);
    \draw[->] (forward) -- (backward);
    \draw[->] (backward) -- (linesearch);
    \draw[->] (linesearch) -- (converge);
    \draw[->] (converge) -- node[align=center, left] {No} (1, 1) -- (forward);
    \draw[->] (converge) -- node[align=center, above right] {Yes} (exit);
  \end{tikzpicture}
  \caption{iLQR algorithm loop: initialize control trajectory, forward rollout
    to compute state trajectory, backward pass to compute feedback gains,
    line search to update controls, and convergence check.}
  \label{fig:ilqr-loop}
\end{figure}

Here we implement iLQR end-to-end:
Differential Dynamic Programming (Volume~I, Chapter~5) is the workhorse for
trajectory optimization. Here we implement it end-to-end:

\begin{lstlisting}[language=Python, caption={iLQR implementation for trajectory optimization.}]
import numpy as np
from typing import Callable, Tuple

def ilqr_solve(
    dynamics: Callable,
    cost_running: Callable,
    cost_terminal: Callable,
    x0: np.ndarray,
    u_init: np.ndarray,
    n_iterations: int = 50,
    regularization: float = 1e-6
) -> Tuple[np.ndarray, np.ndarray, list[float]]:
    """Iterative Linear Quadratic Regulator.

    Parameters
    ----------
    dynamics : callable
        (x, u) -> x_next, (A, B) Jacobians
    cost_running : callable
        (x, u) -> (cost, lx, lu, lxx, luu, lux)
    cost_terminal : callable
        x -> (cost, lx, lxx)
    x0 : np.ndarray
        Initial state.
    u_init : np.ndarray, shape (T, m)
        Initial control trajectory.
    n_iterations : int
        Maximum iterations.
    regularization : float
        Regularization for Quu inversion.

    Returns
    -------
    Tuple
        (x_traj, u_traj, costs)
    """
    T, m = u_init.shape
    n = len(x0)

    # Forward rollout with initial controls
    x_traj = np.zeros((T + 1, n))
    x_traj[0] = x0
    u_traj = u_init.copy()

    for t in range(T):
        x_next, _ = dynamics(x_traj[t], u_traj[t])
        x_traj[t + 1] = x_next

    costs = []

    for iteration in range(n_iterations):
        # Compute total cost
        total_cost = sum(
            cost_running(x_traj[t], u_traj[t])[0]
            for t in range(T)
        ) + cost_terminal(x_traj[T])[0]
        costs.append(total_cost)

        # Backward pass
        Vx_next, Vxx_next = cost_terminal(x_traj[T])[1:]
        k_gains = np.zeros((T, m))
        K_gains = np.zeros((T, m, n))

        for t in range(T - 1, -1, -1):
            _, (A, B) = dynamics(x_traj[t], u_traj[t])
            _, lx, lu, lxx, luu, lux = cost_running(
                x_traj[t], u_traj[t]
            )

            Qx = lx + A.T @ Vx_next
            Qu = lu + B.T @ Vx_next
            Qxx = lxx + A.T @ Vxx_next @ A
            Quu = luu + B.T @ Vxx_next @ B + regularization * np.eye(m)
            Qux = lux + B.T @ Vxx_next @ A

            Quu_inv = np.linalg.inv(Quu)
            k_gains[t] = -Quu_inv @ Qu
            K_gains[t] = -Quu_inv @ Qux

            Vx_next = Qx + K_gains[t].T @ Quu @ k_gains[t]
            Vxx_next = Qxx + K_gains[t].T @ Quu @ K_gains[t]

        # Forward pass with line search
        alpha = 1.0
        x_new = np.zeros_like(x_traj)
        u_new = np.zeros_like(u_traj)
        x_new[0] = x0

        for t in range(T):
            dx = x_new[t] - x_traj[t]
            u_new[t] = u_traj[t] + alpha * k_gains[t] + K_gains[t] @ dx
            x_new[t + 1], _ = dynamics(x_new[t], u_new[t])

        x_traj = x_new
        u_traj = u_new

        if iteration > 0 and abs(costs[-1] - costs[-2]) < 1e-6:
            break

    return x_traj, u_traj, costs
\end{lstlisting}

\section{Application: Optimal Pendulum Swing-Up}

The pendulum swing-up is a classic trajectory optimization benchmark. The goal is to
find the minimum-energy torque sequence that brings the pendulum from the downward
(stable) equilibrium to the upright (unstable) equilibrium.

This is a canonical example in Agrachev \& Sachkov \cite{AgrachevSachkov2004} (Chapter~5, Example 5.1):
the pendulum system is an affine control system $\ddot{\theta} = -g\sin\theta + u$, and the swing-up
problem is optimal synthesis: find the control law that minimizes energy while ensuring convergence
to the upright equilibrium. Their theory predicts that the optimal trajectory will exhibit
``bang-bang'' control (maximum torque in one direction) followed by a final feedback phase---
which is exactly what occurs in practice.
The classic control benchmark: swing a pendulum from the hanging position to
the inverted equilibrium.

\section{Application: Golf Swing Optimization}

Applying DDP to a 3-DOF golf model to optimize clubhead speed at impact.

\section*{Exercises}
\begin{enumerate}
  \item Use the iLQR code to solve the pendulum swing-up problem.
  \item Extend to a double pendulum and optimize for maximum tip speed.
  \item Add muscle force constraints and solve the muscle-driven swing-up.
\end{enumerate}

\chapter{Controller Design and Testing}
\label{ch:controller-design}

\epigraph{A controller is only as good as its worst-case performance.}{}

\section{PD Control with Gravity Compensation}

\begin{lstlisting}[language=Python, caption={PD controller with gravity compensation.}]
import numpy as np

class PDGravityCompController:
    """PD controller with gravity compensation.

    From Volume I, Chapter 3: the simplest nonlinear controller
    that guarantees asymptotic stability.
    """

    def __init__(self, Kp: np.ndarray, Kd: np.ndarray,
                 gravity_func=None):
        self.Kp = Kp
        self.Kd = Kd
        self.gravity_func = gravity_func

    def compute(self, q: np.ndarray, qd: np.ndarray,
                q_des: np.ndarray, qd_des: np.ndarray) -> np.ndarray:
        """Compute control torques.

        Parameters
        ----------
        q, qd : current state
        q_des, qd_des : desired state

        Returns
        -------
        tau : np.ndarray, control torques
        """
        tau = self.Kp @ (q_des - q) + self.Kd @ (qd_des - qd)
        if self.gravity_func is not None:
            tau += self.gravity_func(q)
        return tau
\end{lstlisting}

\section{Trajectory Tracking with Contraction}

Implementing the contraction-based controller from Volume~I, Chapter~4.

\section{Funnel Control}

Implementing trajectory tubes (Volume~II, Chapter~3) for robust tracking.

\section*{Exercises}
\begin{enumerate}
  \item Implement PD control for the 2-DOF arm and track a circular trajectory.
  \item Add gravity compensation and compare tracking error.
  \item Implement a contraction-based controller and verify convergence.
\end{enumerate}

\chapter{Parameter Estimation and Model Fitting}
\label{ch:param-estimation}

\epigraph{All models are wrong, but some are less wrong than others once you fit them.}{}

\section{System Identification for Multi-Body Models}

\begin{lstlisting}[language=Python, caption={Inertial parameter estimation from motion data.}]
import numpy as np
from scipy.optimize import minimize

def estimate_inertial_params(
    q_data: np.ndarray,
    qd_data: np.ndarray,
    qdd_data: np.ndarray,
    tau_data: np.ndarray,
    n_params: int
) -> np.ndarray:
    """Estimate inertial parameters using least-squares.

    The equations of motion are LINEAR in inertial parameters:
    tau = Y(q, qd, qdd) * pi
    where Y is the regressor matrix and pi is the parameter vector.

    Parameters
    ----------
    q_data : (N, n) joint positions
    qd_data : (N, n) joint velocities
    qdd_data : (N, n) joint accelerations
    tau_data : (N, n) joint torques
    n_params : int number of inertial parameters

    Returns
    -------
    np.ndarray estimated parameters
    """
    N, n = q_data.shape

    # Build regressor matrix (simplified for 1-DOF)
    Y = np.zeros((N * n, n_params))
    tau_vec = tau_data.reshape(-1)

    for k in range(N):
        for j in range(n):
            row = k * n + j
            Y[row, 0] = qdd_data[k, j]  # Inertia
            Y[row, 1] = np.sin(q_data[k, j])  # Gravity

    # Least squares solution
    params, _, _, _ = np.linalg.lstsq(Y, tau_vec, rcond=None)
    return params
\end{lstlisting}

\section{Muscle Parameter Calibration}

Using the Hill model from Volume~III with optimization to match experimental data.

\section*{Exercises}
\begin{enumerate}
  \item Estimate the mass and length of a pendulum from simulated motion data.
  \item Add measurement noise and study estimation robustness.
  \item Implement Bayesian estimation with prior knowledge from literature.
\end{enumerate}

\chapter{Reinforcement Learning for Motor Control}
\label{ch:rl}

\epigraph{When the model is too complex for optimal control, let the agent learn.}{}

\section{RL Environments in UpstreamDrift}

UpstreamDrift provides Gymnasium-compatible environments for motor control tasks:

\begin{lstlisting}[language=Python, caption={Custom RL environment for pendulum control.}]
import numpy as np

class PendulumEnv:
    """Simple pendulum environment for RL benchmarks.

    State: [cos(theta), sin(theta), theta_dot]
    Action: torque (continuous, clipped to [-2, 2])
    Reward: -(angle^2 + 0.1*vel^2 + 0.01*action^2)
    """

    def __init__(self, dt: float = 0.05, max_steps: int = 200):
        self.dt = dt
        self.max_steps = max_steps
        self.g = 9.81
        self.l = 1.0
        self.m = 1.0
        self.max_torque = 2.0
        self.state = np.zeros(2)
        self.step_count = 0

    def reset(self, seed: int | None = None) -> np.ndarray:
        if seed is not None:
            np.random.seed(seed)
        self.state = np.array([
            np.pi + np.random.uniform(-0.1, 0.1),
            np.random.uniform(-0.5, 0.5)
        ])
        self.step_count = 0
        return self._get_obs()

    def _get_obs(self) -> np.ndarray:
        theta, theta_dot = self.state
        return np.array([np.cos(theta), np.sin(theta), theta_dot])

    def step(self, action: float) -> tuple:
        theta, theta_dot = self.state
        u = np.clip(action, -self.max_torque, self.max_torque)

        # Dynamics
        theta_ddot = (-self.g / self.l * np.sin(theta) +
                      u / (self.m * self.l**2))
        theta_dot += theta_ddot * self.dt
        theta += theta_dot * self.dt

        self.state = np.array([theta, theta_dot])
        self.step_count += 1

        # Reward
        angle_normalized = ((theta + np.pi) % (2*np.pi)) - np.pi
        reward = -(angle_normalized**2 +
                   0.1 * theta_dot**2 +
                   0.01 * u**2)

        done = self.step_count >= self.max_steps
        return self._get_obs(), reward, done, {}


# Quick test
env = PendulumEnv()
obs = env.reset(seed=42)
total_reward = 0.0
for _ in range(200):
    action = np.random.uniform(-2, 2)
    obs, reward, done, _ = env.step(action)
    total_reward += reward
    if done:
        break
print(f"Random policy total reward: {total_reward:.1f}")
\end{lstlisting}

\section{From Gym to Golf}

The UpstreamDrift RL pipeline extends from simple benchmarks to full golf swing optimization
using MyoSuite muscle-driven environments.

\section*{Exercises}
\begin{enumerate}
  \item Train a PPO agent to swing up the pendulum. Plot the learning curve.
  \item Compare RL-learned policy with the iLQR solution from Chapter~5.
  \item Design a reward function for maximizing clubhead speed in a golf model.
\end{enumerate}

\chapter{Visualization and Analysis Tools}
\label{ch:visualization}

\epigraph{If you cannot see it, you cannot understand it.}{}

\section{3D Motion Visualization}

\begin{lstlisting}[language=Python, caption={Stick-figure animation for multi-body motion.}]
import numpy as np

def compute_joint_positions_2d(
    q: np.ndarray,
    lengths: list[float]
) -> np.ndarray:
    """Compute 2D joint positions for a serial chain.

    Parameters
    ----------
    q : np.ndarray, shape (n,)
        Joint angles.
    lengths : list[float]
        Segment lengths.

    Returns
    -------
    np.ndarray, shape (n+1, 2)
        XY positions of each joint (including base).
    """
    n = len(q)
    positions = np.zeros((n + 1, 2))
    angle_cumulative = 0.0

    for i in range(n):
        angle_cumulative += q[i]
        positions[i + 1, 0] = (positions[i, 0] +
                               lengths[i] * np.cos(angle_cumulative))
        positions[i + 1, 1] = (positions[i, 1] +
                               lengths[i] * np.sin(angle_cumulative))

    return positions


def phase_portrait(
    q_history: np.ndarray,
    qd_history: np.ndarray,
    joint_idx: int = 0,
    label: str = "Joint 0"
) -> dict:
    """Compute phase portrait data for a single joint.

    Parameters
    ----------
    q_history : (N, n) trajectory data
    qd_history : (N, n) velocity data
    joint_idx : int
    label : str

    Returns
    -------
    dict with 'x', 'y', 'label' keys
    """
    return {
        'x': q_history[:, joint_idx],
        'y': qd_history[:, joint_idx],
        'label': label
    }
\end{lstlisting}

\section{Jacobian and Manipulability Visualization}

Visualizing the velocity manipulability ellipsoid $J J^T$ at each configuration.

\section{Energy Flow Diagrams}

Tracking kinetic and potential energy across segments to understand energy transfer
in the kinetic chain.

\section*{Exercises}
\begin{enumerate}
  \item Animate a 3-link arm tracking a circular path.
  \item Plot the manipulability ellipsoid during a reaching movement.
  \item Create an energy flow diagram for the double pendulum golf model.
\end{enumerate}

\chapter{Capstone: The Golf Swing Analysis Pipeline}
\label{ch:golf-capstone}

\begin{gomdraftnotice}
This chapter is under active development and contains only a structural outline
with preliminary material. Full exposition, worked examples, proofs, and exercises
will be added in future editions.
\end{gomdraftnotice}

\epigraph{%
  We have built the tools. Now let us use them all.
}{}

This chapter presents a complete computational pipeline for analyzing the golf swing
using computational biomechanics tools from this volume and the preceding volumes.
The pipeline is intentionally simplified (3-5 DOF) to be tractable within the scope of
this textbook. Real golf swings are complex multi-segment movements with 20+ DOF and
significant muscle redundancy, typically studied via motion capture analysis combined
with musculoskeletal models \cite{siciliano2009}.

\textbf{Important Note on Biomechanical Claims:} All descriptions of swing mechanics in
this chapter are restricted to mathematical derivations from the model or explicit citations
of biomechanical literature. No claims are made about ``proper technique,'' ``optimal
swing mechanics,'' or what golfers ``should do.'' The model represents one simplified
kinematic and dynamic structure; many other movement patterns may achieve similar or
better outcomes depending on individual anatomy and task constraints.

\section{Project Overview}

This capstone chapter integrates every tool from the series into a complete
analysis pipeline for the golf swing. The pipeline implements:

\begin{center}
\renewcommand{\arraystretch}{1.3}
\begin{tabular}{@{}llr@{}}
\toprule
\textbf{Stage} & \textbf{Theory Source} & \textbf{Chapter} \\
\midrule
1. Model Construction & Screw theory, PoE kinematics & Vol~0 \\
2. Forward Dynamics & Lagrangian mechanics & Vol~0 \\
3. Trajectory Optimization & DDP/iLQR & Vol~I \\
4. Contraction Analysis & Contraction theory & Vol~I \\
5. Funnel Computation & Trajectory tubes & Vol~II \\
6. ZTCF family Decomposition & Drift vs.\ control & Vol~I \\
7. Muscle Force Estimation & Hill model + static opt & Vol~III \\
8. Motor Analysis & UCM, synergies & Vol~IV \\
9. Visualization \& Export & Companion platform tools & Vol~V \\
\bottomrule
\end{tabular}
\end{center}

\section{Stage 1: Model}

% DO NOT EDIT. Generated by scripts/generate_worked_examples.py.
% Every number below is solved, not transcribed; CI fails if this file
% drifts from what the generator produces (issue #3518).


\begin{lstlisting}[language=Python, caption={Golf swing analysis pipeline, generated from \texttt{src/affine\_control/swing\_analysis.py}.}]
"""Golf swing analysis pipeline for the Volume V capstone project.

Volume V chapter 10 printed this pipeline as a 160-line code listing captioned
"Complete golf swing analysis pipeline". It was never executed, and three parts
of it did not work:

* ``compute_clubhead_speed`` summed ``qdot_i * r_i`` as scalars, with no
  trigonometry at all. Tip velocity is a *vector* sum whose directions depend on
  the accumulated joint angles, so the printed version was configuration-blind:
  it returned the same speed for a straight arm and a fully folded one. It even
  computed ``cumulative_angle`` and then never used it, and carried a dead
  Jacobian loop whose body was ``pass``.

* ``ztcf_decomposition`` returned a hardcoded 60/40 split of clubhead speed. Its
  own comments said the result "should not be trusted" and that a real
  implementation must solve the forward dynamics twice -- which is what this one
  does.

* The driver assigned ``np.random.randn(N, 3) * 10`` as joint torques, which the
  placeholder decomposition then ignored, so the pipeline was disconnected at
  both ends.

The chapter's listing is now generated from this file, so what the book shows is
what CI runs, and the freshness gate fails if the two drift apart. See #3518.
"""

from __future__ import annotations

from dataclasses import dataclass, field

import numpy as np
from numpy.typing import NDArray

from src.affine_control.golf_model import GolfModel

__all__ = ["SwingAnalysis", "synthetic_swing"]

type Array = NDArray[np.float64]

METRES_PER_SECOND_TO_MPH = 2.236936


def _rod_inertia(mass: float, length: float) -> float:
    """Moment of inertia of a uniform rod about its own centre of mass."""
    return mass * length**2 / 12.0


@dataclass
class SwingAnalysis:
    """Kinematic and dynamic analysis of a three-link planar swing.

    The chain is shoulder, elbow, wrist, carrying the club as the distal link.
    Angles are relative joint angles; the model treats each segment as a uniform
    rod, so its inertia follows from mass and length rather than being a free
    parameter.
    """

    segment_lengths: Array = field(default_factory=lambda: np.array([0.28, 0.46, 1.10]))
    segment_masses: Array = field(default_factory=lambda: np.array([2.16, 1.79, 0.30]))

    time: Array | None = None
    joint_angles: Array | None = None
    joint_velocities: Array | None = None
    joint_torques: Array | None = None

    def model(self) -> GolfModel:
        """The rigid-body model implied by the segment parameters."""
        lengths = tuple(float(v) for v in self.segment_lengths)
        masses = tuple(float(v) for v in self.segment_masses)
        inertias = tuple(_rod_inertia(m, ell) for m, ell in zip(masses, lengths, strict=True))
        return GolfModel(masses=masses, lengths=lengths, inertias=inertias)  # type: ignore[arg-type]

    def _require_kinematics(self) -> tuple[Array, Array, Array]:
        """Return ``(time, angles, rates)``, or fail with a usable message."""
        if self.time is None or self.joint_angles is None or self.joint_velocities is None:
            raise ValueError("set time, joint_angles and joint_velocities before analysing")
        return self.time, self.joint_angles, self.joint_velocities

    def clubhead_speed(self) -> Array:
        """Clubhead speed at each sample, from the tip Jacobian.

        The tip velocity is ``sum_i qdot_i * (z_hat x r_i)``, where ``r_i`` runs
        from joint ``i`` to the tip. Its magnitude depends on the configuration
        through the angles between those contributions -- which is precisely what
        a scalar sum of ``qdot_i * |r_i|`` throws away.
        """
        _, angles, rates = self._require_kinematics()
        model = self.model()
        return np.array([model.clubhead_speed(q, qd) for q, qd in zip(angles, rates, strict=True)])

    def ztcf_decomposition(self) -> tuple[Array, Array]:
        """Split clubhead acceleration into drift and control contributions.

        Solves the forward dynamics twice at each sample -- once with the applied
        torques and once with zero torque -- and attributes the difference to
        control. Because the dynamics are control-affine, that difference is
        exactly ``M^-1 tau`` and the split is exact rather than approximate.

        Returns the joint-space acceleration magnitudes ``(drift, control)``.
        """
        _, angles, rates = self._require_kinematics()
        if self.joint_torques is None:
            raise ValueError("set joint_torques before decomposing")
        model = self.model()

        drift = np.zeros(len(angles))
        control = np.zeros(len(angles))
        for index, (q, qd, tau) in enumerate(zip(angles, rates, self.joint_torques, strict=True)):
            bias = model.coriolis(q, qd) @ qd + model.gravity_torque(q)
            mass = model.rigid_mass_matrix(q)
            drift_accel = np.linalg.solve(mass, -bias)
            control_accel = np.linalg.solve(mass, tau)
            drift[index] = float(np.linalg.norm(drift_accel))
            control[index] = float(np.linalg.norm(control_accel))
        return drift, control

    def summary(self) -> str:
        """Human-readable summary of the analysis."""
        time, _, _ = self._require_kinematics()
        speeds = self.clubhead_speed()
        peak = float(np.max(speeds))
        peak_index = int(np.argmax(speeds))

        lines = [
            "=== Golf Swing Analysis Summary ===",
            f"Duration: {time[-1]:.3f} s",
            f"Peak clubhead speed: {peak:.1f} m/s " f"({peak * METRES_PER_SECOND_TO_MPH:.1f} mph)",
            f"Peak speed time: {time[peak_index]:.3f} s",
        ]

        if self.joint_torques is not None:
            drift, control = self.ztcf_decomposition()
            # Report the median over interior samples rather than a value at one
            # instant. Two reasons: the endpoints carry one-sided np.gradient
            # estimates of the joint rates, and the drift term legitimately
            # passes through zero whenever the chain straightens -- at a fully
            # extended vertical configuration both the gravity and Coriolis
            # torques vanish, so a ratio quoted there reads as "0% drift" and
            # says nothing about the swing.
            interior = slice(1, -1)
            total = drift[interior] + control[interior]
            share = float(np.median(100.0 * drift[interior] / total))
            lines.append(f"Median over the swing: drift {share:.0f}%, control {100 - share:.0f}%")
        return "\n".join(lines)


def synthetic_swing(duration: float = 0.30, dt: float = 0.001) -> SwingAnalysis:
    """A smooth synthetic swing, for demonstrating the pipeline.

    These are not measured data. Real joint torques must come from inverse
    dynamics on motion capture, or from a Hill-type muscle model driven by EMG;
    the torques here are a smooth profile chosen only so the decomposition has
    something non-trivial to separate. The printed listing used
    ``np.random.randn``, which made the output different on every run and, since
    the placeholder decomposition ignored the torques entirely, changed nothing.
    """
    samples = int(duration / dt)
    time = np.linspace(0.0, duration, samples)
    fraction = time / duration

    angles = np.column_stack(
        [
            -np.pi / 2 + np.pi * fraction**2,
            np.pi / 3 * np.sin(np.pi * fraction),
            -np.pi / 2 * np.cos(np.pi * fraction / 2),
        ]
    )
    rates = np.gradient(angles, dt, axis=0)
    torques = np.column_stack(
        [
            120.0 * (1.0 - fraction),
            60.0 * np.sin(np.pi * fraction),
            25.0 * fraction**2,
        ]
    )
    return SwingAnalysis(
        time=time, joint_angles=angles, joint_velocities=rates, joint_torques=torques
    )
\end{lstlisting}


\begin{warningbox}
The listing above is generated from \texttt{src/affine\_control/swing\_analysis.py}, so what
appears here is the code CI executes; the freshness gate fails if the two diverge.

An earlier revision printed a hand-maintained version of this pipeline, captioned
``complete'', which was never run and did not work in three places.

\texttt{compute\_clubhead\_speed} summed $\dot q_i r_i$ as scalars, with no trigonometry.
Tip velocity is a \emph{vector} sum whose contributions point in different directions
depending on the joint angles, so the printed version was configuration-blind: it returned
the same speed for a straight arm and a folded one. Substituting a folded configuration into
both gives $13.3$~m/s from the correct Jacobian against $22.5$~m/s from the scalar sum ---
a $69\%$ overestimate. The old code even accumulated \texttt{cumulative\_angle} and then
never used it, and carried a Jacobian loop whose body was \texttt{pass}.

\texttt{ZTCF\_decomposition} returned a hardcoded $60/40$ split of clubhead speed. Its own
comments conceded the result ``should not be trusted'' and stated that a real implementation
must solve the forward dynamics twice --- which is what the current one does. Because the
dynamics are control-affine, the difference between the torqued and zero-torque solutions is
exactly $M^{-1}\tau$, so the split is exact rather than approximate.

The driver assigned \texttt{np.random.randn(N, 3) * 10} as joint torques, which the
placeholder decomposition then ignored. The pipeline was disconnected at both ends.
\end{warningbox}

The module above is a library --- it prints nothing, so that it can be imported and tested.
To exercise it:

\begin{lstlisting}[language=Python, caption={Running the pipeline.}]
from src.affine_control.swing_analysis import synthetic_swing

print(synthetic_swing().summary())
\end{lstlisting}

\begin{remark}[Reading the summary output]
This reports a peak clubhead speed near $30$~m/s for the synthetic kinematics.
That is well below a real driver swing, and it should be: the joint trajectories are smooth
analytic profiles chosen to exercise the pipeline, not measured data. Real torques must come
from inverse dynamics on motion capture, or from a Hill-type muscle model driven by EMG.

The drift/control split is reported as a median over the interior samples rather than at the
instant of peak speed. The endpoints carry one-sided \texttt{np.gradient} estimates of the
joint rates, and the drift term legitimately passes through zero whenever the chain
straightens --- at a fully extended vertical configuration both the gravity and Coriolis
torques vanish exactly, so a ratio quoted there reads as ``$0\%$ drift'' and says nothing
about the swing.
\end{remark}

\section{Golf Swing Simulation Architecture}

\begin{figure}[h]
  \centering
  \begin{tikzpicture}[
    every node/.style={align=center, draw, rounded corners, minimum width=1.8cm, text centered, font=\tiny},
    ]
    % Row 1: Input
    \node[fill=blue!20] (model) at (0, 6) {Model:\\Kinematics};
    \node[fill=blue!20] (motion) at (2, 6) {Motion:\\Data};
    \node[fill=blue!20] (forces) at (4, 6) {Force\\Plates};

    % Row 2: Dynamics
    \node[fill=green!20] (forward) at (1, 4.5) {Forward\\Dynamics};
    \node[fill=green!20] (inverse) at (3, 4.5) {Inverse\\Dynamics};

    % Row 3: Analysis
    \node[fill=yellow!20] (ZTCF) at (0, 3) {ZTCF};
    \node[fill=yellow!20] (muscle) at (2, 3) {Muscle\\Force};
    \node[fill=yellow!20] (ucm) at (4, 3) {UCM};

    % Row 4: Output
    \node[fill=orange!20] (report) at (2, 1.5) {Summary\\Report};

    % Edges
    \draw[->] (model) -- (forward);
    \draw[->] (motion) -- (inverse);
    \draw[->] (motion) -- (forward);
    \draw[->] (forces) -- (inverse);
    \draw[->] (forward) -- (ZTCF);
    \draw[->] (inverse) -- (muscle);
    \draw[->] (forward) -- (ucm);
    \draw[->] (ZTCF) -- (report);
    \draw[->] (muscle) -- (report);
    \draw[->] (ucm) -- (report);
  \end{tikzpicture}
  \caption{Complete golf swing analysis architecture. Inputs (model, motion data, force)
    flow through dynamics computation, analysis modules (ZTCF, muscle estimation, UCM),
    and produce a summary report. Each box references the corresponding chapter from Volumes 0--IV.}
  \label{fig:golf-arch}
\end{figure}

\section{Stage 2--9: Full Pipeline}

Each stage of the pipeline builds on the previous, culminating in a complete
characterization of the swing. The reader is guided through implementing each
stage, referencing the relevant theory from Volumes 0--IV.

\textbf{Data Requirements:} A complete pipeline requires:
\begin{itemize}
  \item Motion capture data: joint angles $\bm{q}(t)$ sampled at sufficient frequency (120+ Hz typical).
  \item Force measurements: ground reaction forces and moments (optional but recommended).
  \item Anatomical model: segment lengths, masses, and inertias (from literature or imaging).
\end{itemize}

\section*{Final Exercises}

\begin{enumerate}
  \item \textbf{Complete Pipeline.} Run the analysis pipeline on the synthetic swing.
        Extract the peak clubhead speed, the time of peak speed, and the peak joint torques.
        Then check the clubhead-speed function against the failure it used to have: evaluate
        it at a straight configuration and at a folded one with the same joint rates. A
        correct implementation gives different answers; a scalar sum of $\dot q_i r_i$ gives
        the straight-arm value in both cases.

  \item \textbf{Parameter Study.} Programmatically vary club length ($\pm 10\%$) and mass
        ($\pm 20\%$). Simulate the swing and measure peak clubhead speed. Plot clubhead
        speed versus each parameter.

  \item \textbf{Muscle Analysis.} Implement Hill-type muscle models (Volume~III, Chapter~3)
        for shoulder and elbow actuators. Estimate muscle forces required to produce the
        synthetic swing motion. Which muscle groups have the highest peak force?

  \item \textbf{Robustness Analysis.} Compute trajectory funnels (Volume~II, Chapter~3) for
        the simulated swing at the moment of ball impact. Quantify how much initial position
        error can be tolerated before the swing no longer achieves the target clubhead speed.

  \item \textbf{Variability Analysis.} Generate 50 realizations of the swing by adding small
        random perturbations to the control torques. Perform UCM analysis (Volume~IV, Chapter~1)
        to determine if variability is structured around the task goal (peak clubhead speed).

  \item \textbf{Open Project.} Choose a different multi-joint movement (e.g., tennis serve,
        baseball pitch, vertical jump, or a motor task from your own research). Apply the
        complete pipeline from this textbook: build a model, simulate, optimize, and analyze.
        Document in a Jupyter notebook with plots.
\end{enumerate}


\backmatter
\printindex

\bibliographystyle{plain}
\bibliography{../geometry_of_motion}

\end{document}
